\section{The standard control and its first-order arithmetic}
\label{m:sec:control}
Define, literally,
\begin{equation}\label{m:eq:f}
 F_0(n)=2^n,\qquad F_{j+1}(n)=F_j^{\,n+1}(n),\qquad
 f(n)=F_5(n+100).
\end{equation}
A superscript on a function denotes a finite iterate. The starting argument
in $F_{j+1}$ varies with $n$. This convention will not be replaced by
iteration from a fixed seed.

\begin{lemma}[Growth]\label{m:lem:growth}
Each $F_j$ is recursive and strictly increasing, with $F_j(n)\ge2^n$.
For $j\ge0$,
\[
 F_{j+1}(n+1)\ge2^{F_{j+1}(n)+1}.
\]
Consequently $f(z+1)\ge(A+2)f(z)$ whenever $z\ge A\ge0$.
\end{lemma}
\begin{proof}
Finite iteration preserves recursiveness, strict increase, and domination
of $2^n$. Increasing the input from $n$ to $n+1$ makes the first $n+1$
iterates at least one larger, and the definition then supplies one further
application of $F_j$. Thus
$F_{j+1}(n+1)\ge F_j(F_{j+1}(n)+1)\ge2^{F_{j+1}(n)+1}$.
For $v\ge1$, $2^{v+1}\ge(v+2)v$, by direct checking at $v=1,2,3$ and
induction. Apply this with $v=f(z)>z\ge A$.
\end{proof}

\begin{lemma}[Effective constant residues]\label{m:lem:residues}
For every $m\ge1$ one can compute $r_m\in[0,m)$ such that
\[
 F_j(n)\equiv r_m\pmod m\quad(j\ge1,\ n\ge m).
\]
In particular $f(n)\equiv r_m\pmod m$ for $n\ge\max(0,m-100)$.
For $s\ge1$, $j\ge2$, this can be strengthened to
\begin{equation}\label{m:eq:envelope}
 m\le 2^{2^s}\quad\Longrightarrow\quad F_j(s)\equiv r_m\pmod m.
\end{equation}
If $B=f(u)$ and $m\le2^{2^B}$, then also $f(u+1)\equiv r_m\pmod m$.
\end{lemma}
\begin{proof}
Set $r_1=0$. For $m>1$, write $m=2^a b$ with $b$ odd, and put
$q=\varphi(b)$ if $b>1$, $q=1$ otherwise. Since $q<m$, recursively
compute $r_q$, then use the Chinese remainder theorem to define
\[
 r_m\equiv0\pmod{2^a},\qquad r_m\equiv2^{r_q}\pmod b.
\]
The recursion strictly lowers the integer modulus.
For towers $T_0(n)=n$, $T_{h+1}(n)=2^{T_h(n)}$, induction on $m$ shows
$T_h(n)\equiv r_m\pmod m$ whenever $h,n\ge m$. Indeed, the last exponent
is at least $n\ge a$, while its residue modulo $q$ is $r_q$ by induction;
Euler's theorem handles the odd factor. This proves the assertion for
$F_1(n)=T_{n+1}(n)$. At every higher level, all intermediate arguments are
at least $n$, so the last application of the preceding level proves the
same residue assertion.

For the stronger bound, $F_1(s)\ge2^{2^s}$. In $F_2(s)$ there are at
least two $F_1$ applications; the final one has input at least $F_1(s)$,
so its input is at least $m$. The same argument propagates through every
$j\ge2$. Finally put $n=u+100$. The last $F_4$ call in
$F_5(n+1)=F_4^{n+2}(n+1)$ has input
$F_4^{n+1}(n+1)>F_4^{n+1}(n)=B$. Apply the stronger bound to that call.
\end{proof}

\begin{theorem}[Snyatkov's arithmetic supplier]\label{m:thm:snyatkov}
The first-order theory $\operatorname{Th}(\N;0,1,+,<,f)$, with the exact
function in \eqref{m:eq:f}, is decidable.
\end{theorem}
\begin{proof}
We use Theorem~2 and effective Corollary~2 of Snyatkov~\cite{snyatkov},
pp.~50--51. Its definitions on p.~40 require an increasing recursive
function effectively compatible with addition: effective eventual periodicity
modulo every integer, and an effective shifted sign separation for each
unbounded integer combination of fixed translates. It also requires effective
eventual periodicity of the \emph{fixed-seed} hyperfunction
$H(0)=1$, $H(n+1)=f(H(n))$. The cited theorem and effective corollary
give effective quantifier elimination after adjoining the specified
divisibility and floor-inverse symbols for $f$ and $H$.  Both functions
are recursive, strictly increasing and unbounded.  Their floor inverses
are computable by finite search, with value zero below the initial
function value.  The remaining atomic predicates are decidable.
Quantifier elimination therefore decides closed sentences in the
expanded language and in its $f$-only reduct.

The residue condition follows from Lemma~\ref{m:lem:residues}. For the second
condition combine equal shifts in $S(x)=\sum_b a_b f(x+b)$. If every
combined coefficient is zero the sum is zero. Otherwise let $B$ be the
largest shift with nonzero coefficient, and set
\[
 A=\sum_{b<B}|a_b|,\qquad
 \delta=\max(0,-\min b,A+1-B).
\]
Every argument at $x+\delta$ is nonnegative. With
$z=x+\delta+B-1\ge x+A$, Lemma~\ref{m:lem:growth} gives, if $a_B>0$,
\[
 S(x+\delta)\ge f(z+1)-Af(z)\ge2f(z)>f(x).
\]
For $a_B<0$, apply the same calculation to $-S$. This proves unboundedness,
its sign, and a suitable shift by finite arithmetic on the input data.
Lastly $H(n)\ge n+1$; for $n\ge m$, the last application of $f$ in
$H(n)$ has input at least $m$, so $H(n)\equiv r_m\pmod m$.
All published hypotheses are therefore paid. No identification of $H$ with
a diagonal $F_j$, and no decision result for variable-seed iteration, is
being imported.
\end{proof}

